\documentclass[11pt,fontset=fandol]{ctexart}
\usepackage[a4paper,margin=23mm]{geometry}
\usepackage{amsmath,amssymb,mathtools,bm,mathrsfs,esint}
\usepackage{graphicx,longtable,booktabs,array,calc}
\usepackage[normalem]{ulem}
\usepackage{xcolor,hyperref}
\hypersetup{hidelinks}
\setcounter{secnumdepth}{-1}
\newcounter{none}
\providecommand{\tightlist}{\setlength{\itemsep}{0pt}\setlength{\parskip}{0pt}}
\providecommand{\passthrough}[1]{#1}
\usepackage{amssymb,mathrsfs,esint,bm,newunicodechar}
\allowdisplaybreaks
\setlength{\emergencystretch}{3em}
\providecommand{\llangle}{\langle\!\langle}
\providecommand{\rrangle}{\rangle\!\rangle}
\newunicodechar{∞}{\ensuremath{\infty}}
\newunicodechar{∈}{\ensuremath{\in}}
\newunicodechar{∑}{\ensuremath{\sum}}
\newunicodechar{∫}{\ensuremath{\int}}
\newunicodechar{≤}{\ensuremath{\leq}}
\newunicodechar{≥}{\ensuremath{\geq}}
\newunicodechar{≠}{\ensuremath{\neq}}
\newunicodechar{→}{\ensuremath{\to}}
\newunicodechar{⇒}{\ensuremath{\Rightarrow}}
\newunicodechar{α}{\ensuremath{\alpha}}
\newunicodechar{β}{\ensuremath{\beta}}
\newunicodechar{γ}{\ensuremath{\gamma}}
\newunicodechar{λ}{\ensuremath{\lambda}}
\newunicodechar{φ}{\ensuremath{\varphi}}
\newunicodechar{ε}{\ensuremath{\varepsilon}}
\newunicodechar{σ}{\ensuremath{\sigma}}
\newunicodechar{θ}{\ensuremath{\theta}}
\newunicodechar{δ}{\ensuremath{\delta}}
\newunicodechar{η}{\ensuremath{\eta}}
\newunicodechar{π}{\ensuremath{\pi}}
\newunicodechar{ξ}{\ensuremath{\xi}}
\newunicodechar{Σ}{\ensuremath{\Sigma}}
\newunicodechar{Δ}{\ensuremath{\Delta}}
\newunicodechar{∂}{\ensuremath{\partial}}
\newunicodechar{∀}{\ensuremath{\forall}}
\newunicodechar{∃}{\ensuremath{\exists}}
\newunicodechar{√}{\ensuremath{\surd}}
\newunicodechar{−}{\ensuremath{-}}
\newunicodechar{μ}{\ensuremath{\mu}}
\newunicodechar{ρ}{\ensuremath{\rho}}
\newunicodechar{τ}{\ensuremath{\tau}}
\newunicodechar{ω}{\ensuremath{\omega}}
\newunicodechar{Ω}{\ensuremath{\Omega}}
\newunicodechar{χ}{\ensuremath{\chi}}
\newunicodechar{ψ}{\ensuremath{\psi}}
\newunicodechar{Γ}{\ensuremath{\Gamma}}
\newunicodechar{Φ}{\ensuremath{\Phi}}
\newunicodechar{Λ}{\ensuremath{\Lambda}}
\newunicodechar{ℝ}{\ensuremath{\mathbb{R}}}
\newunicodechar{ℤ}{\ensuremath{\mathbb{Z}}}
\newunicodechar{ℕ}{\ensuremath{\mathbb{N}}}
\newunicodechar{△}{\ensuremath{\triangle}}
\newunicodechar{⊗}{\ensuremath{\otimes}}
\begin{document}
\section{2024年杭州电子科技大学601数学分析真题及答案(答案不全版)}\label{ux5e74ux676dux5ddeux7535ux5b50ux79d1ux6280ux5927ux5b66601ux6570ux5b66ux5206ux6790ux771fux9898ux53caux7b54ux6848ux7b54ux6848ux4e0dux5168ux7248}

来源：\texttt{HDU/601数学分析/2024/2024年杭州电子科技大学601数学分析真题及答案(答案不全版).pdf}。原
PDF 页序：1, 2。

\begin{quote}
本文件保留原``答案不全版''所含的同年题面，依据同年回忆版交叉校正公式记号；重复答案页已在清单中映射至其他版本。
\end{quote}

\section{杭州电子科技大学2024年数学分析（回忆版）真题}\label{ux676dux5ddeux7535ux5b50ux79d1ux6280ux5927ux5b662024ux5e74ux6570ux5b66ux5206ux6790ux56deux5fc6ux7248ux771fux9898}

\begin{quote}
文字整理说明：根据远程OCR与本地逐页OCR整理数学语法，保留回忆版缺失内容；四2第二个函数的两份OCR识读有分歧，已单独标记。
\end{quote}

\subsection{一、填空题（8题，每题5分，共40分）}\label{ux4e00ux586bux7a7aux98988ux9898ux6bcfux98985ux5206ux517140ux5206}

\begin{enumerate}
\def\labelenumi{\arabic{enumi}.}
\item
  \(\displaystyle\lim_{n\to\infty}(n^{2024}+2^n)^{1/n}=\) \_\_\_\_。
\item
  \(\displaystyle\int\frac{x^3}{\sqrt{1-x^2}}\,dx=\) \_\_\_\_。
\item
  \(y=2^x\) 的麦克劳林级数中 \(x^{2024}\) 的系数为 \_\_\_\_。
\item
  极限函数 \(\displaystyle\lim_{n\to\infty}\frac{(n+1)\sin x}{nx+4}\)
  的间断点为 \_\_\_\_。
\item
  求 \(x^2+y^2+z^2=6\) 与 \(x^2+y^2=z\) 的交线的切线方程。
\item
  没回忆到。
\item
  \(f(x,y,z)=x^2+xy+\cdots\)，求从 \(P_0(-1,0,2)\) 指向 \(P_1(2,1,1)\)
  的方向导数。（函数表达式未回忆全；\(P_1\)
  坐标由原解析题干与本地OCR共同恢复。）
\item
  \(L\) 为 \(x^2+y^2=1\) 的上半圆周，则
  \(\displaystyle\int_L(x^2+y^2)\,ds=\) \_\_\_\_。
\end{enumerate}

\subsection{二、计算题（共25分）}\label{ux4e8cux8ba1ux7b97ux9898ux517125ux5206}

\begin{enumerate}
\def\labelenumi{\arabic{enumi}.}
\item
  （5分）\(\displaystyle\lim_{x\to0}\frac{x-\sin x}{x(e^x-1)}\)。
\item
  （10分）求
  \(\displaystyle\sum_{n=1}^\infty(-1)^n\frac{2n+1}{n}x^{2n}\)
  的收敛域与和函数。
\item
  （10分）求
  \(\displaystyle\int_L(-e^y\cos x-xy)\,dx-e^y\sin x\,dy\)，其中 \(L\)
  为 \(x^2+y^2=x\) 从 \((1,0)\) 到 \((0,0)\) 的上半圆周。
\end{enumerate}

\subsection{三、证明题（每题10分，共30分）}\label{ux4e09ux8bc1ux660eux9898ux6bcfux989810ux5206ux517130ux5206}

\begin{enumerate}
\def\labelenumi{\arabic{enumi}.}
\item
  \(f(x,y)=|x-y|\varphi(x,y)\)，\(\varphi\) 在 \((0,0)\) 连续。 （1）求
  \(f_x(0,0),f_y(0,0)\) 存在的充要条件； （2）讨论 \(f\) 在 \((0,0)\)
  的可微性。
\item
  \(\displaystyle\iint_\Sigma(\cdot)\,dy\,dz+(\cdot)\,dz\,dx+(\cdot)\,dx\,dy\)，\(\Sigma:z=5-x^2-y^2\)
  在 \(x\ge1\)
  的部分，方向取上侧。（原资料未回忆到具体积分系数；其他边界条件也需核实。）
\item
  含参量反常积分
  \(\displaystyle I(x)=\int_0^\infty e^{-t}\frac{\sin(tx)}t\,dt\)，求
  \(I(x)\)。
\end{enumerate}

\subsection{四、综合题（共55分）}\label{ux56dbux7efcux5408ux9898ux517155ux5206}

\begin{enumerate}
\def\labelenumi{\arabic{enumi}.}
\item
  （10分）已知
  \(\displaystyle\lim_{x\to x_0^+}f(x)=A\)，用``\(\varepsilon\)---\(\delta\)''语言证明：\(\displaystyle\lim_{x\to x_0^+}\sqrt[3]{f(x)}=\sqrt[3]A\)。
\item
  （10分）证明 \(f(x)=x+\sin^2x\) 在 \(\mathbb R\)
  上一致连续，\(g(x)=x+\sin(x^2)\) 在 \(\mathbb R\)
  上不一致连续。（第二个函数按远程OCR的 \(x+\sin x^2\)
  整理；本地OCR识读为 \(x+\sin^2x\)，存在分歧，见补充解析。）
\item
  （10分）\(f\) 二阶可导，\(f(0)=0,f(1)=1,f'(0)=f'(1)=0\)。证明存在
  \(\xi\in(0,1)\) 使 \(|f''(\xi)|\ge4\)。
\item
  （10分）\(f\) 有一阶连续导数，\(f(0)=f(1)=0\)。证明 \[
  \left|\int_0^1f(x)\,dx\right|\le\frac14\max_{x\in[0,1]}|f'(x)|.
  \]
\item
  （15分）设函数项级数
  \(\displaystyle f(x)=\sum_{n=1}^\infty\frac{e^{-nx}}{n^2+1}\)。证明：
  （1）\(f\) 在 \([0,\infty)\) 上连续； （2）\(f\) 在 \((0,\infty)\)
  上可导； （3）\(f\) 在 \(x=0\) 不可导。
\end{enumerate}
\end{document}
